\clearpage

\section{Original sizing and annual results}\label{app:d}

\subsection*{Alternative put sizing pipelines retained from the original study}

\begin{table}[H]\centering\small
\caption{Alternative put-sizing policies on the common reporting split.}
\setlength{\tabcolsep}{4pt}\renewcommand{\arraystretch}{1.18}
\begin{tabularx}{\textwidth}{@{}>{\hsize=1.75781\hsize\linewidth=\hsize\raggedright\arraybackslash}X>{\hsize=0.78125\hsize\linewidth=\hsize\raggedleft\arraybackslash}X>{\hsize=0.78125\hsize\linewidth=\hsize\raggedleft\arraybackslash}X>{\hsize=0.70312\hsize\linewidth=\hsize\raggedleft\arraybackslash}X>{\hsize=0.97656\hsize\linewidth=\hsize\raggedleft\arraybackslash}X@{}}\toprule
\textbf{Put sizing policy} & \textbf{2018--21 S} & \textbf{2022--25 S} & \textbf{Later entries} & \textbf{Later profit}\\\midrule
Constant 50\% & 0.437 & 1.007 & 151 & \$0.205m\\
Constant 75\% & 0.449 & 0.944 & 104 & \$0.195m\\
Constant 100\% & 0.478 & 1.055 & 75 & \$0.201m\\
Tree + CEM & 0.573 & 1.263 & 110 & \$0.249m\\
Spread-rank sizing & 0.464 & 0.815 & 121 & \$0.165m\\
Ridge-score sizing & 0.421 & 0.890 & 94 & \$0.180m\\
Direct linear policy & 0.267 & 0.873 & 96 & \$0.171m\\
Direct additive policy & 0.302 & 1.013 & 97 & \$0.180m\\
\bottomrule\end{tabularx}\end{table}

All policies receive the same put proposals, directions, contracts and funded admission rules. The lower four rows are follow-up comparisons rather than original selection candidates. Tree/CEM has the highest funded Sharpe in both periods among these policies. These comparisons retain the original question about entry sizing alongside the newer straddle and allocation experiments.


\begin{table}[H]\centering\small
\caption{Put-bandit annual account returns and excess performance.}
\setlength{\tabcolsep}{4pt}\renewcommand{\arraystretch}{1.18}
\begin{tabularx}{\textwidth}{@{}>{\hsize=0.58594\hsize\linewidth=\hsize\raggedright\arraybackslash}X>{\hsize=1.07422\hsize\linewidth=\hsize\raggedleft\arraybackslash}X>{\hsize=1.22070\hsize\linewidth=\hsize\raggedleft\arraybackslash}X>{\hsize=1.17188\hsize\linewidth=\hsize\raggedleft\arraybackslash}X>{\hsize=0.94727\hsize\linewidth=\hsize\raggedleft\arraybackslash}X@{}}\toprule
\textbf{Year} & \textbf{Account return} & \textbf{Above cash (pp)} & \textbf{Excess Sharpe} & \textbf{Entries}\\\midrule
2022 & 3.470\% & 1.673 & 2.458 & 66\\
2023 & 5.288\% & 0.432 & 1.492 & 19\\
2024 & 5.104\% & -0.283 & -0.907 & 11\\
2025* & 3.207\% & 0.169 & 1.092 & 14\\
\bottomrule\end{tabularx}\end{table}

Calendar-period account returns and above-cash differences are not annualized, including partial 2025. Sharpe is annualized. In 2024 the account grows while losing relative to cash. The original put strategy's results are concentrated in a few pairs and in 2022; this limits claims of broad repeated success.


The original put results are concentrated: SNAP/UBER, GAP/TJX and AMD/MU contribute 66.9\% of later net pair trading profit, and 2022 supplies 83.1\%. These figures exclude shared financing. They describe where gains occurred, not a rule for selecting those names retrospectively. Annual Sharpe ratios are retained in Appendix~\ref{app:c}.


The put account's later annualized total account return is 4.667\%, versus 0.530\% growth relative to cash. Average reserve is 2.44\% and peak reserve 15.77\%. Its cash-relative drawdown is 0.495\%, compared with 0.646\% at full size. These figures retain the original distinction between a profitable trading signal, account funding and capital deployment.


\FloatBarrier
